% open-logic.tex
% compiles to produce complete text using open-logic-dev style

\documentclass[../include/open-logic-part]{subfiles}

\begin{document}

\clearpage

\begin{editorial}
આ ફાઇલ ઓપન લોજિક પ્રોજેક્ટની બધી સામગ્રી લોડ કરે છે.
આવી સંપાદકીય નોંધો દેખાય, તો તેનો અર્થ એ છે કે સામગ્રી
કેવી રીતે રજૂ થશે તે વિચાર્યા વિના ફાઇલ સંકલિત થઈ છે.
આ નોંધ વાંચી શકો, તો કદાચ આ પીડીએફમાંથી ભણાવવું કે
અભ્યાસ કરવો \emph{યોગ્ય નથી}.

ઓપન લોજિક પ્રોજેક્ટ ભણાવવા કે સ્વાધ્યાય માટે વધુ અનુકૂળ
પાઠ બનાવવા અનેક વ્યવસ્થાઓ આપે છે.
દાખલા તરીકે, મૂળભૂત ગોઠવણીમાં બધા તાર્કિક સંયોજકો મૂળભૂત ગણાય છે,
અને સાબિતીઓમાં દરેક સંયોજકના બધા કિસ્સા પૂરા કરવામાં આવે છે.
પરંતુ કેટલાક કિસ્સા કસરત તરીકે છોડવા વધુ સારું છે.
ઓપન લોજિક પ્રોજેક્ટનું કાર્ય પણ હજી ચાલુ છે.
સહયોગ અને સુધારો પ્રોત્સાહિત કરવા સામગ્રી ફક્ત મસૌદામાં હોય,
કસરતો અધૂરી હોય, વગેરે સ્થિતિમાં પણ તેને સમાવાય છે.
અભ્યાસક્રમ માટે તૈયાર થતી પીડીએફમાં આવા વિભાગો બાકાત રહે.

ભણાવવા અને અભ્યાસ માટે વધુ યોગ્ય પીડીએફ શોધવા
\href{http://builds.openlogicproject.org/}{પ્રોજેક્ટના જાળસ્થાને ઉપલબ્ધ નમૂના અભ્યાસક્રમો} જુઓ.
પોતાનો પાઠ બનાવવા શરૂઆત
\href{https://github.com/OpenLogicProject/OpenLogic/tree/master/courses/sample}{નમૂના મુખ્ય ફાઇલ}થી કરી શકાય;
વધુ સુગઠિત અને આગળનાં ઉદાહરણો માટે આ પ્રોજેક્ટ પરથી
બનેલાં પાઠ્યપુસ્તકોના સ્રોત જુઓ.
\end{editorial}

\olimport[sets-functions-relations]{sets-functions-relations-complete}

\olimport[propositional-logic]{propositional-logic}

\olimport[first-order-logic]{first-order-logic}

\olimport[model-theory]{model-theory}

\olimport[computability]{computability}

\olimport[turing-machines]{turing-machines}

\olimport[incompleteness]{incompleteness}

\olimport[second-order-logic]{second-order-logic}

\olimport[lambda-calculus]{lambda-calculus}

\olimport[many-valued-logic]{many-valued-logic}

\olimport[normal-modal-logic]{normal-modal-logic}

\olimport[applied-modal-logic]{applied-modal-logic}

\olimport[intuitionistic-logic]{intuitionistic-logic}

\olimport[counterfactuals]{counterfactuals}

\olimport[set-theory]{set-theory}

\olimport[methods]{methods}

\olimport[history]{history}

\olimport[reference]{reference}

\end{document}

